\documentclass[conference]{IEEEtran}
\IEEEoverridecommandlockouts
% The preceding line is only needed to identify funding in the first footnote. If that is unneeded, please comment it out.
%Template version as of 6/27/2024

\usepackage{cite}
\usepackage{amsmath,amssymb,amsfonts}
\usepackage{algorithmic}
\usepackage{graphicx}
\usepackage{textcomp}
\usepackage{xcolor}
\usepackage{tabularx}
\usepackage{array}

\newcolumntype{Y}{>{\raggedright\arraybackslash}X}

\def\BibTeX{{\rm B\kern-.05em{\sc i\kern-.025em b}\kern-.08em
    T\kern-.1667em\lower.7ex\hbox{E}\kern-.125emX}}
\begin{document}

\title{Reliable Extraction from Epilepsy Letters\\}

\author{\IEEEauthorblockN{1\textsuperscript{st} Conor Brown}
\IEEEauthorblockA{\textit{dept. name of organization (of Aff.)} \\
\textit{name of organization (of Aff.)}\\
Belfast, Northern Ireland \\
cbrown564@gmail.com}
}

\maketitle

\begin{abstract}
This document is a model and instructions for \LaTeX.
This and the IEEEtran.cls file define the components of your paper [title, text, heads, etc.]. *CRITICAL: Do Not Use Symbols, Special Characters, Footnotes, 
or Math in Paper Title or Abstract.
\end{abstract}

\begin{IEEEkeywords}
component, formatting, style, styling, insert.
\end{IEEEkeywords}

\section{Introduction}
\subsection{Medical framing}
Epilepsy is one of the most common neurological diseases

30 percent of patients are drug resistant and research progress has been stubbornly slow

We need a deeper understanding of who these patients are, how they are treated, how outcomes differ over time and across groups

The most important clinical markers are trapped in unstructured notes

Costly chart review is required for constructing patient timelines, identifying research cohorts, conducting retrospective studies

They struggle through that anyway because it’s so valuable, but it’s not just costly it’s error prone 
Entity extraction vs. entity relational extraction and reasoning
NLP solved entity extraction a while back, but this is a temporal and clinical reasoning problem - how to represent a clinical state in a rigorous way

It’s not just computers that struggle with this problem, but doctors too - they all describe things differently and experts struggle to interpret them identically 

\subsection{Prior system summary}
Rules-based systems are good at entity extraction but require extensive annotation and perform poorly on logic

BERT-based systems require extensive annotations and again, perform poorly on logic

GPT-based systems are very semantically flexible and strong on reasoning, but they’re a bit too flexible - hallucination, rambling prose, it’s unstructured data to unstructured data

Monolithic LLMs provide the least visibility into the most dangerous task, clinical reasoning

Most of the LLMs tested and performing well as monolithic systems were closed-source, which creates privacy concerns

These problems lead to systems generalising poorly to new tasks and styles
We need the modularity and reliability of rules-based systems to constrain the flexibility of the LLMs and audit their outputs

\subsection{Hybrid solution}
I present a modular, evidence-grounded LLM framework for clinical information extraction from epilepsy clinical letters

It decomposes extraction into configurable stages for semantic interpretation, relation extraction, temporal resolution, concept normalisation, schema validation and evidence verification

It is evaluated across two independently developed, publicly shareable epilepsy benchmarks assessing two different capabilities: broad phenotype extraction in ExECTv2 and temporally complex seizure-frequency extraction from synthetic leets

By directly comparing rules-only, LLM-only and hybrid configurations, plus component ablations and error analysis, we show how modular, evidence-grounded hybrid systems not only improve field-level F1 scores but also reliability, auditability and generalisability

\section{Figures and Tables}\label{FAT}
\paragraph{Positioning Figures and Tables} Place figures and tables at the top and 
bottom of columns. Avoid placing them in the middle of columns. Large 
figures and tables may span across both columns. Figure captions should be 
below the figures; table heads should appear above the tables. Insert 
figures and tables after they are cited in the text. Use the abbreviation 
``Fig.~\ref{fig}'', even at the beginning of a sentence.

\begin{tabularx}{\linewidth}{|Y|Y|Y|Y|}
\hline
\textbf{Table Head} &
\textbf{Claim it supports} &
\textbf{Evidence shown} &
\textbf{Reviewer risk} \\
\hline

Table 1: Dataset Summary &
Method is robust and generalisable &
Dataset size, labels, splits, class balance &
How valuable is synthetic data? \\
\hline

Table 2: Main Results &
Method improves over baseline in 3 key ways &
F1, evidence, schema validity, uncertainty &
What do schema and evidence validity mean? \\
\hline

Figure 1: Architecture &
Method is modular, transparent, reproducible, and clinically logical &
Pipeline stages, representative example &
Architecture may seem unnecessarily convoluted \\
\hline

Table 3: Ablation &
Deterministic and LLM components are complementary &
LLM-only, rules-only, LLM + schema validation, LLM + evidence verifier,
LLM + concept normalisation, LLM + temporal resolver, full hybrid &
Unclear separation between rules and LLM; unclear rules \\
\hline

Figure 2: Error Analysis &
Deterministic and LLM components have distinct error profiles that you can calibrate for &
Hallucinations, low recall, schema validity, evidence coverage,
normalisation errors, temporal contradiction, abstention rate &
Difficulty identifying the most meaningful errors \\
\hline
\end{tabularx}

\begin{figure}[htbp]
\centerline{\includegraphics{fig1.png}}
\caption{Example of a figure caption.}
\label{fig}
\end{figure}

Figure Labels: Use 8 point Times New Roman for Figure labels. Use words 
rather than symbols or abbreviations when writing Figure axis labels to 
avoid confusing the reader. As an example, write the quantity 
``Magnetization'', or ``Magnetization, M'', not just ``M''. If including 
units in the label, present them within parentheses. Do not label axes only 
with units. In the example, write ``Magnetization (A/m)'' or ``Magnetization 
\{A[m(1)]\}'', not just ``A/m''. Do not label axes with a ratio of 
quantities and units. For example, write ``Temperature (K)'', not 
``Temperature/K''.

\section*{Acknowledgment}

The preferred spelling of the word ``acknowledgment'' in America is without 
an ``e'' after the ``g''. Avoid the stilted expression ``one of us (R. B. 
G.) thanks $\ldots$''. Instead, try ``R. B. G. thanks$\ldots$''. Put sponsor 
acknowledgments in the unnumbered footnote on the first page.

\section*{References}

Please number citations consecutively within brackets \cite{b1}. The 
sentence punctuation follows the bracket \cite{b2}. 

\begin{thebibliography}{00}
\bibitem{b1} G. Eason, B. Noble, and I. N. Sneddon, ``On certain integrals of Lipschitz-Hankel type involving products of Bessel functions,'' Phil. Trans. Roy. Soc. London, vol. A247, pp. 529--551, April 1955.
\bibitem{b2} J. Clerk Maxwell, A Treatise on Electricity and Magnetism, 3rd ed., vol. 2. Oxford: Clarendon, 1892, pp.68--73.
\bibitem{b3} I. S. Jacobs and C. P. Bean, ``Fine particles, thin films and exchange anisotropy,'' in Magnetism, vol. III, G. T. Rado and H. Suhl, Eds. New York: Academic, 1963, pp. 271--350.
\bibitem{b4} K. Elissa, ``Title of paper if known,'' unpublished.
\bibitem{b5} R. Nicole, ``Title of paper with only first word capitalized,'' J. Name Stand. Abbrev., in press.
\bibitem{b6} Y. Yorozu, M. Hirano, K. Oka, and Y. Tagawa, ``Electron spectroscopy studies on magneto-optical media and plastic substrate interface,'' IEEE Transl. J. Magn. Japan, vol. 2, pp. 740--741, August 1987 [Digests 9th Annual Conf. Magnetics Japan, p. 301, 1982].
\bibitem{b7} M. Young, The Technical Writer's Handbook. Mill Valley, CA: University Science, 1989.
\bibitem{b8} D. P. Kingma and M. Welling, ``Auto-encoding variational Bayes,'' 2013, arXiv:1312.6114. [Online]. Available: https://arxiv.org/abs/1312.6114
\bibitem{b9} S. Liu, ``Wi-Fi Energy Detection Testbed (12MTC),'' 2023, gitHub repository. [Online]. Available: https://github.com/liustone99/Wi-Fi-Energy-Detection-Testbed-12MTC
\bibitem{b10} ``Treatment episode data set: discharges (TEDS-D): concatenated, 2006 to 2009.'' U.S. Department of Health and Human Services, Substance Abuse and Mental Health Services Administration, Office of Applied Studies, August, 2013, DOI:10.3886/ICPSR30122.v2
\bibitem{b11} K. Eves and J. Valasek, ``Adaptive control for singularly perturbed systems examples,'' Code Ocean, Aug. 2023. [Online]. Available: https://codeocean.com/capsule/4989235/tree
\end{thebibliography}

\vspace{12pt}
\color{red}
IEEE conference templates contain guidance text for composing and formatting conference papers. Please ensure that all template text is removed from your conference paper prior to submission to the conference. Failure to remove the template text from your paper may result in your paper not being published.

\end{document}
